\label{part:mathematical-formalism}

\begin{introduction}
    \item \textbf{Part II}: Mathematical Formalism - Fiber Bundles \& Tensor Fields
    \item 从哲学思辨到严格数学的惊险一跃
    \item 构成方程：$\Psi = \mathbf{T}_S \otimes \mathbf{V}_Q$
    \item 纤维丛：上帝收纳宇宙的文件夹
\end{introduction}

这一部分完成了从哲学思辨到严格数学的惊险一跃。我们引入了\textbf{纤维丛}，这是上帝收纳宇宙的文件夹，也是我们理解复杂性的终极几何工具。

通过推导\textbf{构成方程 ($\Psi = \mathbf{T}_S \otimes \mathbf{V}_Q$)}，我们将“创造”这一神圣的行为，还原为了一次精确的\textbf{张量积运算}。我们证明了，实体（Entity）不过是质在形的流形上切出的一个\textbf{截面}，而变化（Change）不过是联络引导下的\textbf{平行移动}。

现在，我们拥有了描述万物的数学语言。但这套语言真的能描述我们所处的这个物理世界吗？它仅仅是柏拉图天堂里的幻梦，还是现实世界的投影？为了回答这个问题，我们需要拿起这把数学的尺子，去丈量真实的物理学。

%--chp---
